% problem_id: S046-thm:ttac
% source_id: S046 (HoTT Book)
% source_file: basics.tex
% statement_locator: basics.tex lines 2418--2421
% proof_locator: basics.tex lines 2422--2436
% env: thm
% pairing: tight (verified)
% status: complete (statement + author proof verbatim + reason + locators)
%
\begin{thm}\label{thm:ttac}
  \index{universal!property!of dependent pair type}%
  \eqref{eq:sigma-ump-map} is an equivalence.
\end{thm}

\begin{proof}
  As before, we define a quasi-inverse to send $(g,h)$ to the function $\lam{x} (g(x),h(x))$.
  Now given $f:\prd{x:X} \sm{a:A(x)} P(x,a)$, the round-trip composite yields the function
  \begin{equation}
    \lam{x} (\proj1(f(x)),\proj2(f(x))).\label{eq:prod-ump-rt2}
  \end{equation}
  Now for any $x:X$, by \cref{thm:eta-sigma} (the uniqueness principle for $\Sigma$-types) we have
  %
  \begin{equation*}
    (\proj1(f(x)),\proj2(f(x))) = f(x).
  \end{equation*}
  %
  Thus, by function extensionality,~\eqref{eq:prod-ump-rt2} is equal to $f$.
  On the other hand, given $(g,h)$, the round-trip composite yields $(\lam {x} g(x),\lam{x} h(x))$, which is judgmentally equal to $(g,h)$ as before.
\end{proof}
